\documentclass[12pt, a4paper]{article}

\usepackage[utf8]{inputenc}
\usepackage[T1]{fontenc}
\usepackage[french]{babel}
\usepackage[margin=2.5cm]{geometry}
\usepackage{xcolor}
\usepackage{booktabs}
\usepackage{longtable}
\usepackage{array}
\usepackage{enumitem}
\usepackage{titlesec}
\usepackage{fancyhdr}
\usepackage{mdframed}

% ── Couleurs ──────────────────────────────────────────────────────────────────
\definecolor{dxcpurple}{RGB}{102, 0, 153}
\definecolor{blockbg}{RGB}{248, 242, 255}
\definecolor{greenbg}{RGB}{232, 245, 233}
\definecolor{orangebg}{RGB}{255, 243, 224}
\definecolor{greytext}{RGB}{90, 90, 90}

% ── En-tête / pied de page ────────────────────────────────────────────────────
\pagestyle{fancy}
\fancyhf{}
\fancyhead[L]{\textcolor{dxcpurple}{\textbf{DXC Technology}}}
\fancyhead[R]{\textcolor{dxcpurple}{Cahier des Charges — Tableau de Bord Sprint}}
\fancyfoot[C]{\thepage}
\renewcommand{\headrulewidth}{0.4pt}
\renewcommand{\headrule}{%
  \hbox to\headwidth{\color{dxcpurple}\leaders\hrule height \headrulewidth\hfill}}

% ── Titres ────────────────────────────────────────────────────────────────────
\titleformat{\section}{\large\bfseries\color{dxcpurple}}{}{0em}{}[\titlerule]
\titleformat{\subsection}{\normalsize\bfseries\color{dxcpurple}}{}{0em}{}

% ── Environnements encadrés ───────────────────────────────────────────────────
\newmdenv[
  backgroundcolor=blockbg,
  linecolor=dxcpurple,
  linewidth=1.2pt,
  innerleftmargin=12pt, innerrightmargin=12pt,
  innertopmargin=10pt, innerbottommargin=10pt,
  roundcorner=4pt
]{purplebox}

\newmdenv[
  backgroundcolor=greenbg,
  linecolor=green!60!black,
  linewidth=1pt,
  innerleftmargin=12pt, innerrightmargin=12pt,
  innertopmargin=8pt, innerbottommargin=8pt,
  roundcorner=4pt
]{greenbox}

\newmdenv[
  backgroundcolor=orangebg,
  linecolor=orange!80!black,
  linewidth=1pt,
  innerleftmargin=12pt, innerrightmargin=12pt,
  innertopmargin=8pt, innerbottommargin=8pt,
  roundcorner=4pt
]{orangebox}

% ─────────────────────────────────────────────────────────────────────────────
\begin{document}

% ── Page de titre ─────────────────────────────────────────────────────────────
\begin{center}
  {\color{dxcpurple}\rule{\linewidth}{1.5pt}}\\[0.5cm]
  {\LARGE\bfseries\color{dxcpurple}
    Tableau de Bord de Pilotage de Sprint}\\[0.25cm]
  {\large IRPAUTO · DXC GraphTalk Service Desk}\\[0.15cm]
  {\normalsize\color{greytext}
    Cahier des charges et spécifications fonctionnelles}\\[0.4cm]
  {\color{dxcpurple}\rule{\linewidth}{1.5pt}}
\end{center}

\vspace{0.6cm}

% ─────────────────────────────────────────────────────────────────────────────
\section{Contexte Métier}

DXC Technology est un prestataire mondial de services informatiques qui accompagne ses
clients dans la modernisation et la maintenance de leurs systèmes d'information. Dans
le cadre du projet \textbf{IRPAUTO}, l'équipe prend en charge le maintien en condition
opérationnelle du progiciel \textbf{GraphTalk} pour le compte de plusieurs compagnies
d'assurance françaises, sur des domaines métier variés~: prévoyance, santé, comptabilité
et édition documentaire.

\vspace{0.2cm}

Le travail est cadencé en \textbf{sprints Agile} d'environ quatre semaines. À chaque
sprint, l'équipe s'engage à livrer un ensemble de tickets — corrections de bugs,
évolutions fonctionnelles, incidents de production — préalablement sélectionnés et
estimés. Cet engagement constitue le \textbf{périmètre du sprint}, et sa bonne
exécution est un indicateur clé de la relation de confiance avec le client.

\vspace{0.4cm}

% ─────────────────────────────────────────────────────────────────────────────
\section{Problématique}

\begin{orangebox}
Le suivi de l'avancement du sprint reposait jusqu'alors sur des
\textbf{processus manuels} dont les limites devenaient de plus en plus
contraignantes~:

\begin{itemize}[leftmargin=1.5em, itemsep=3pt]
  \item \textbf{Absence de vision en temps réel}~: les données n'étaient
        disponibles qu'à travers des fichiers produits ponctuellement, rapidement
        obsolètes dès qu'un ticket changeait d'état.
  \item \textbf{Aucune vue consolidée}~: il n'existait pas de tableau de bord
        centralisé permettant de visualiser l'avancement global, ni par équipe
        ni par individu.
  \item \textbf{Pas de suivi de trajectoire}~: il était impossible de savoir,
        à un instant donné, si l'équipe était en avance ou en retard sur le
        rythme nécessaire pour honorer ses engagements à la fin du sprint.
  \item \textbf{Reporting chronophage}~: la préparation des comptes rendus
        hebdomadaires nécessitait des consolidations manuelles longues et
        sujettes aux erreurs.
  \item \textbf{Manque de transparence}~: les parties prenantes (management,
        client) ne disposaient pas d'un accès simple et actualisé à l'état
        réel du sprint.
\end{itemize}
\end{orangebox}

\vspace{0.4cm}

% ─────────────────────────────────────────────────────────────────────────────
\section{Objectifs du Projet}

\begin{purplebox}
\textbf{Objectif général~:} concevoir et déployer un outil de pilotage qui
fournit, à tout moment, une vision claire, fiable et consolidée de l'avancement
du sprint — sans intervention manuelle — aussi bien pour l'équipe interne que
pour les présentations client.
\end{purplebox}

\vspace{0.4cm}

\subsection{Objectifs fonctionnels}

\begin{enumerate}[leftmargin=2em, itemsep=6pt]

  \item \textbf{Automatiser la collecte des données.}\\
        Supprimer la dépendance aux exports manuels. Les données doivent être
        récupérées automatiquement depuis la source officielle de gestion des
        tickets, de façon périodique ou à la demande.

  \item \textbf{Fournir une vue d'avancement journalière du sprint (burndown).}\\
        L'outil doit permettre de visualiser, pour chaque jour du sprint~:
        le volume de travail restant, le volume livré, et la trajectoire
        théorique idéale. Cette courbe permet de détecter tout écart par rapport
        aux engagements pris.

  \item \textbf{Mesurer le taux d'engagement par individu et par équipe.}\\
        Pour chaque ticket du périmètre du sprint, le système doit déterminer
        s'il est considéré comme \emph{livré} ou \emph{non livré} sur la base
        de son statut courant. Cette classification alimente des indicateurs
        d'engagement (tickets livrés, Story Points livrés, taux d'engagement en
        pourcentage) à la maille de chaque collaborateur et de chaque squad.

  \item \textbf{Suivre l'évolution du backlog des incidents de production.}\\
        En parallèle du sprint, l'outil doit permettre de surveiller le volume
        global d'incidents ouverts, d'en observer la tendance quotidienne
        (entrées, sorties, solde net) et de le comparer à une cible de réduction.

  \item \textbf{Produire des rapports visuels prêts pour les présentations client.}\\
        Les données calculées doivent être disponibles dans un outil de
        visualisation externe (type PowerBI) permettant de générer des rapports
        graphiques professionnels sans manipulation supplémentaire.

\end{enumerate}

\vspace{0.4cm}

% ─────────────────────────────────────────────────────────────────────────────
\section{Spécifications Fonctionnelles}

\subsection{Module 1 — Burndown du Sprint}

\begin{itemize}[leftmargin=1.5em, itemsep=3pt]
  \item Affichage jour par jour de l'avancement du sprint sur toute sa durée.
  \item Indicateurs calculés~: points restants, points résolus, trajectoire idéale,
        pourcentage de complétion.
  \item Même logique disponible en nombre de tickets (en plus des Story Points).
  \item Prise en charge d'une logique de date de complétion robuste~: la date
        retenue pour un ticket est la première date disponible parmi la date de
        qualification, la date de résolution et la date de clôture.
  \item Possibilité de configurer manuellement les dates de début et de fin de sprint.
\end{itemize}

\vspace{0.2cm}

\subsection{Module 2 — Engagement de l'Équipe}

\begin{itemize}[leftmargin=1.5em, itemsep=3pt]
  \item Périmètre restreint aux tickets officiellement engagés dans le sprint
        (scope validé lors de la réunion de planification).
  \item Classification de chaque ticket selon son statut courant~:
  \begin{itemize}
    \item \textbf{Livré}~: statuts indiquant qu'un ticket a atteint une étape
          de validation ou de clôture.
    \item \textbf{Non livré}~: statuts indiquant un travail encore en cours.
  \end{itemize}
  \item Calcul par collaborateur~: tickets engagés, SP engagés, tickets livrés,
        SP livrés, tickets non livrés, SP non livrés, taux d'engagement.
  \item Agrégation identique par squad.
  \item Formule du taux d'engagement~:
        $\text{Commitment} = \dfrac{\text{SP livrés}}{\text{SP engagés}} \times 100$
  \item Vue drill-down~: par squad, puis par individu, avec la liste détaillée de
        chaque ticket et son statut courant.
\end{itemize}

\vspace{0.2cm}

\subsection{Module 3 — Backlog des Incidents}

\begin{itemize}[leftmargin=1.5em, itemsep=3pt]
  \item Suivi de l'évolution quotidienne du nombre d'incidents de production ouverts.
  \item Indicateurs~: incidents ouverts, incidents fermés, solde net journalier,
        courbe de tendance, cible de réduction.
  \item Visualisation de l'écart entre la trajectoire réelle et la trajectoire cible.
\end{itemize}

\vspace{0.2cm}

\subsection{Module 4 — Reporting Externe}

\begin{itemize}[leftmargin=1.5em, itemsep=3pt]
  \item Les données calculées (burndown, engagement par squad, engagement par
        individu, liste de tickets classifiés) doivent être exportées dans un
        format structuré accessible à un outil de visualisation externe.
  \item Les rapports doivent se mettre à jour sans intervention manuelle
        (simple rafraîchissement).
  \item Les visuels cibles sont~: courbes de burndown, graphiques d'engagement
        empilés, tableaux de synthèse, matrices squad / individu, drill-through
        par ticket.
\end{itemize}

\vspace{0.4cm}

% ─────────────────────────────────────────────────────────────────────────────
\section{Exigences Non Fonctionnelles}

\begin{longtable}{>{\bfseries\small\color{dxcpurple}}p{4cm}
                  >{\small}p{10cm}}
  \toprule
  \normalfont\textbf{\color{dxcpurple}Critère} &
  \normalfont\textbf{\color{dxcpurple}Exigence} \\
  \midrule
  \endhead
  \bottomrule
  \endfoot

  Actualisation &
  Les données doivent pouvoir être mises à jour à la demande ou de façon
  planifiée, sans intervention manuelle sur les fichiers sources. \\[4pt]

  Fiabilité &
  Source unique de vérité~; les données affichées doivent être cohérentes entre
  tous les modules et avec la réalité du gestionnaire de tickets. \\[4pt]

  Accessibilité &
  Le tableau de bord interne doit être accessible depuis un navigateur web, sans
  installation logicielle pour les utilisateurs. \\[4pt]

  Maintenabilité &
  L'ajout d'un nouveau sprint ou d'un nouveau périmètre ne doit pas nécessiter
  de modifications structurelles de l'outil. \\[4pt]

  Confidentialité &
  Les données (noms, tickets, capacités) sont à usage interne~; l'accès à
  l'application et à la base de données doit être protégé. \\

\end{longtable}

\vspace{0.4cm}

% ─────────────────────────────────────────────────────────────────────────────
\section{Bénéfices Attendus}

\begin{greenbox}
\begin{itemize}[leftmargin=1.5em, itemsep=4pt]
  \item \textbf{Gain de temps}~: élimination des consolidations manuelles~;
        les indicateurs sont disponibles en temps quasi-réel.
  \item \textbf{Meilleure prise de décision}~: la visibilité sur la trajectoire
        du sprint permet d'anticiper les risques et d'ajuster les priorités avant
        qu'il ne soit trop tard.
  \item \textbf{Transparence envers le client}~: un reporting visuel clair et
        actualisé renforce la confiance et facilite les réunions de suivi.
  \item \textbf{Responsabilisation de l'équipe}~: chaque collaborateur peut
        consulter son propre taux d'engagement et l'état de ses tickets en
        autonomie.
  \item \textbf{Réplicabilité}~: la structure de l'outil est conçue pour
        s'appliquer aux sprints futurs sans refonte.
\end{itemize}
\end{greenbox}

\vspace{0.5cm}
\noindent{\color{dxcpurple}\rule{\linewidth}{0.8pt}}
\begin{center}
  \small\color{greytext}
  Document interne — DXC Technology · Projet IRPAUTO
\end{center}

\end{document}
